% GENERATED FILE. Edit canonical lessons, then run npm run generate:books.
% canonical-source-hash: fnv1a64:11a166afcf62b2b7
% canonical-lessons: KA-C374-jadu-magic, KA-C374-pavada, KA-C374-husi, KA-C374-saddu, KA-C374-kollu

\chapter{Magic and Miracles}
\label{ch:ka-magic-and-miracles}
\hlchaptermodality{374}

\begin{chapteropening}
\textbf{By the end of this chapter:} \emph{I can say magic, a miracle, false, a noise and to kill.}

Complete the last lesson of chapter 374: I can say magic, a miracle, false, a noise and to kill.
\end{chapteropening}

\section[\texorpdfstring{\kn{ಜಾದೂ} (jādū)}{ಜಾದೂ (jādū)}]{\kn{ಜಾದೂ} (jādū) — magic; a magic trick}

\label{lesson:KA-C374-jadu-magic}



\begin{quote}
Before the new one: say the Kannada for a narrative, then the Kannada for a triumph.
\end{quote}

\subsection*{You'll want to know: \kn{ಜಾದೂ}}
\textbf{\kn{ಜಾದೂ}} — \emph{jādū} — \textquotedblleft{}magic; a magic trick\textquotedblright{}.

\begin{quote}
\kn{ಜಾತ್ರೆಯಲ್ಲಿ} \kn{ಒಬ್ಬ} \kn{ಮನುಷ್ಯ} \kn{ಮಕ್ಕಳಿಗೆ} \kn{ಜಾದೂ} \kn{ತೋರಿಸಿದನು}.
\end{quote}

\begin{quote}
\emph{jātreyalli obba manuṣya makkaḷige jādū tōrisidanu.}
\end{quote}

\begin{quote}
At the fair a man showed the children magic tricks.
\end{quote}

A magician on stage is a \kn{ಜಾದೂಗಾರ}.

\begin{cousinweb}
\textbf{\kn{ಜಾದೂ}} (\emph{jādū}) is a loan from Persian \emph{jādū}, \textquotedblleft{}magic, sorcery\textquotedblright{}, which came through Hindustani.
\end{cousinweb}

\begin{grammarlens}[title={What you've built}]
One more word for telling what happened.
\end{grammarlens}

\subsection*{Your turn}
\begin{itemize}
\raggedright
  \item \emph{Say it:} \emph{jādū}
  \item \emph{Say it:} \emph{jādū}, once more
  \item \emph{Say it:} \emph{jātreyalli obba manuṣya makkaḷige jādū tōrisidanu.} — At the fair a man showed the children magic tricks.
  \item \emph{Recall:} say the Kannada for a narrative, then the Kannada for a triumph, then say \emph{jādū} again
\end{itemize}

\begin{tcolorbox}[breakable,colback=teal!4,colframe=teal!35!black,title={Before you move on}]
What does \kn{ಜಾದೂ} mean? (\textquotedblleft{}Magic; a magic trick\textquotedblright{}.) Say it once more.
\end{tcolorbox}
\section[\texorpdfstring{\kn{ಪವಾಡ} (pavāḍa)}{ಪವಾಡ (pavāḍa)}]{\kn{ಪವಾಡ} (pavāḍa) — a miracle; a wonder}

\label{lesson:KA-C374-pavada}



\begin{quote}
Before the new one: say the Kannada for a triumph, then the Kannada for magic.
\end{quote}

\subsection*{You'll want to know: \kn{ಪವಾಡ}}
\textbf{\kn{ಪವಾಡ}} — \emph{pavāḍa} — \textquotedblleft{}a miracle; a wonder\textquotedblright{}.

\begin{quote}
\kn{ಆ} \kn{ಅಪಘಾತದಲ್ಲಿ} \kn{ಮಗು} \kn{ಬದುಕುಳಿದಿದ್ದು} \kn{ಒಂದು} \kn{ಪವಾಡ}.
\end{quote}

\begin{quote}
\emph{ā apaghātadalli magu badukuḷididdu ondu pavāḍa.}
\end{quote}

\begin{quote}
It was a miracle that the child survived that accident.
\end{quote}

Plural: \kn{ಪವಾಡಗಳು}, \emph{pavāḍagaḷu}.

\kn{ಪವಾಡ} \kn{ಮಾಡು}, to work a miracle.

\begin{cousinweb}
\textbf{\kn{ಪವಾಡ}} (\emph{pavāḍa}) is a Kannada word whose origin is not settled.
\end{cousinweb}

\begin{grammarlens}[title={What you've built}]
One more word for telling what happened.
\end{grammarlens}

\subsection*{Your turn}
\begin{itemize}
\raggedright
  \item \emph{Say it:} \emph{pavāḍa}
  \item \emph{Say it:} \emph{pavāḍa}, once more
  \item \emph{Say it:} \emph{ā apaghātadalli magu badukuḷididdu ondu pavāḍa.} — It was a miracle that the child survived that accident.
  \item \emph{Recall:} say the Kannada for a triumph, then the Kannada for magic, then say \emph{pavāḍa} again
\end{itemize}

\begin{tcolorbox}[breakable,colback=teal!4,colframe=teal!35!black,title={Before you move on}]
What does \kn{ಪವಾಡ} mean? (\textquotedblleft{}A miracle; a wonder\textquotedblright{}.) Say it once more.
\end{tcolorbox}
\section[\texorpdfstring{\kn{ಹುಸಿ} (husi)}{ಹುಸಿ (husi)}]{\kn{ಹುಸಿ} (husi) — false; empty}

\label{lesson:KA-C374-husi}



\begin{quote}
Before the new one: say the Kannada for magic, then the Kannada for a miracle.
\end{quote}

\subsection*{You'll want to know: \kn{ಹುಸಿ}}
\textbf{\kn{ಹುಸಿ}} — \emph{husi} — \textquotedblleft{}false; empty\textquotedblright{}.

\begin{quote}
\kn{ಆ} \kn{ಜಾದೂಗಾರನ} \kn{ಮಾತು} \kn{ಹುಸಿ} \kn{ಆಯಿತು}.
\end{quote}

\begin{quote}
\emph{ā jādūgārana mātu husi āyitu.}
\end{quote}

\begin{quote}
The magician's words turned out to be false.
\end{quote}

\kn{ಹುಸಿ} \kn{ಬೆದರಿಕೆ}, an empty threat. \kn{ಹುಸಿ} is also a lie.

\begin{cousinweb}
\textbf{\kn{ಹುಸಿ}} (\emph{husi}) is native Dravidian: DEDR 4531 lists Kannada \emph{pusi}, \textquotedblleft{}to become or prove untrue, to lie\textquotedblright{}, and as a noun \textquotedblleft{}a falsehood\textquotedblright{}; old \emph{p-} became \emph{h-}.
\end{cousinweb}

\begin{grammarlens}[title={What you've built}]
One more word for telling what happened.
\end{grammarlens}

\subsection*{Your turn}
\begin{itemize}
\raggedright
  \item \emph{Say it:} \emph{husi}
  \item \emph{Say it:} \emph{husi}, once more
  \item \emph{Say it:} \emph{ā jādūgārana mātu husi āyitu.} — The magician's words turned out to be false.
  \item \emph{Recall:} say the Kannada for magic, then the Kannada for a miracle, then say \emph{husi} again
\end{itemize}

\begin{tcolorbox}[breakable,colback=teal!4,colframe=teal!35!black,title={Before you move on}]
What does \kn{ಹುಸಿ} mean? (\textquotedblleft{}False; empty\textquotedblright{}.) Say it once more.
\end{tcolorbox}
\section[\texorpdfstring{\kn{ಸದ್ದು} (saddu)}{ಸದ್ದು (saddu)}]{\kn{ಸದ್ದು} (saddu) — a noise; a sound}

\label{lesson:KA-C374-saddu}



\begin{quote}
Before the new one: say the Kannada for a miracle, then the Kannada for false.
\end{quote}

\subsection*{You'll want to know: \kn{ಸದ್ದು}}
\textbf{\kn{ಸದ್ದು}} — \emph{saddu} — \textquotedblleft{}a noise; a sound\textquotedblright{}.

\begin{quote}
\kn{ರಾತ್ರಿ} \kn{ಅಡುಗೆಮನೆಯಲ್ಲಿ} \kn{ಏನೋ} \kn{ಸದ್ದು} \kn{ಆಯಿತು}.
\end{quote}

\begin{quote}
\emph{rātri aḍugemaneyalli ēnō saddu āyitu.}
\end{quote}

\begin{quote}
At night there was some noise in the kitchen.
\end{quote}

\kn{ಸದ್ದು} \kn{ಮಾಡು}, to make a noise. \kn{ಸದ್ದಿಲ್ಲದೆ}, without a sound.

\begin{cousinweb}
\textbf{\kn{ಸದ್ದು}} (\emph{saddu}) is usually traced to Sanskrit \emph{śabda}, \textquotedblleft{}sound, noise\textquotedblright{} (CDIAL 12298), through a Prakrit form \emph{sadda}.
\end{cousinweb}

\begin{grammarlens}[title={What you've built}]
One more word for telling what happened.
\end{grammarlens}

\subsection*{Your turn}
\begin{itemize}
\raggedright
  \item \emph{Say it:} \emph{saddu}
  \item \emph{Say it:} \emph{saddu}, once more
  \item \emph{Say it:} \emph{rātri aḍugemaneyalli ēnō saddu āyitu.} — At night there was some noise in the kitchen.
  \item \emph{Recall:} say the Kannada for a miracle, then the Kannada for false, then say \emph{saddu} again
\end{itemize}

\begin{tcolorbox}[breakable,colback=teal!4,colframe=teal!35!black,title={Before you move on}]
What does \kn{ಸದ್ದು} mean? (\textquotedblleft{}A noise; a sound\textquotedblright{}.) Say it once more.
\end{tcolorbox}
\section[\texorpdfstring{\kn{ಕೊಲ್ಲು} (kollu)}{ಕೊಲ್ಲು (kollu)}]{\kn{ಕೊಲ್ಲು} (kollu) — to kill}

\label{lesson:KA-C374-kollu}



\begin{quote}
Before the new one: say the Kannada for false, then the Kannada for a noise.
\end{quote}

\subsection*{You'll want to know: \kn{ಕೊಲ್ಲು}}
\textbf{\kn{ಕೊಲ್ಲು}} — \emph{kollu} — \textquotedblleft{}to kill\textquotedblright{}.

\begin{quote}
\kn{ಕಥೆಯಲ್ಲಿ} \kn{ವೀರ} \kn{ರಾಜಕುಮಾರ} \kn{ಆ} \kn{ರಾಕ್ಷಸನನ್ನು} \kn{ಕೊಂದನು}.
\end{quote}

\begin{quote}
\emph{katheyalli vīra rājakumāra ā rākṣasanannu kondanu.}
\end{quote}

\begin{quote}
In the story the brave prince killed the demon.
\end{quote}

Past (I …): \kn{ಕೊಂದೆ}, \emph{konde}.

The past is irregular: \kn{ಕೊಂದೆ}, I killed; \kn{ಕೊಂದನು}, he killed.

\begin{cousinweb}
\textbf{\kn{ಕೊಲ್ಲು}} (\emph{kollu}) is native Dravidian: DEDR 2132 lists Kannada \emph{kol}, \emph{kollu}, \textquotedblleft{}to kill\textquotedblright{}, and \emph{kole}, \textquotedblleft{}a killing\textquotedblright{}.
\end{cousinweb}

\begin{grammarlens}[title={What you've built}]
One more word for telling what happened.
\end{grammarlens}

\subsection*{Your turn}
\begin{itemize}
\raggedright
  \item \emph{Say it:} \emph{kollu}
  \item \emph{Say it:} \emph{kollu}, once more
  \item \emph{Say it:} \emph{katheyalli vīra rājakumāra ā rākṣasanannu kondanu.} — In the story the brave prince killed the demon.
  \item \emph{Recall:} say the Kannada for false, then the Kannada for a noise, then say \emph{kollu} again
\end{itemize}

\begin{tcolorbox}[breakable,colback=teal!4,colframe=teal!35!black,title={Before you move on}]
What does \kn{ಕೊಲ್ಲು} mean? (\textquotedblleft{}To kill\textquotedblright{}.) Say it once more.
\end{tcolorbox}
